\documentclass[11pt,a4paper]{article}

\usepackage[margin=2.6cm]{geometry}
\usepackage{amsmath,amssymb}
\usepackage{booktabs}
\usepackage{enumitem}
\usepackage{graphicx}
\usepackage{pgfplots}
\pgfplotsset{compat=1.18}
\usepgfplotslibrary{groupplots}
\usepackage[hidelinks]{hyperref}

\newcommand{\dt}{\Delta t}
% A lattice snapshot, framed so that empty ground is distinguishable from paper.
\newcommand{\lattice}[1]{%
  \setlength{\fboxsep}{0pt}\setlength{\fboxrule}{0.4pt}%
  \fbox{\includegraphics[width=\dimexpr\linewidth-0.8pt\relax]{figures/#1}}}
\newcommand{\panelcap}[1]{\\[3pt]{\small #1}}
\newcommand{\code}[1]{\texttt{#1}}

\title{The Multiplicative Dissipation Cellular Automaton}
\author{A specification of the model as implemented}
\date{}

\begin{document}
\maketitle

\begin{abstract}
\noindent
A stochastic cellular automaton on a square lattice in which occupied cells belong to
a two-level taxonomy of \emph{species} and \emph{sub-species}. Cells compete pairwise for
territory, and the outcome of a contest is biased by a species-level energy variable
$\alpha$. Energy is produced in proportion to a species' area and dissipated in proportion
to its internal diversity, so that a species which accumulates too many sub-species
exhausts its energy and \emph{fragments}: its sub-species are promoted, all at once, to
independent species. Speciation is therefore not an elementary move but a catastrophe
that follows a slow build-up of internal multiplicity --- the ``multiplicative
dissipation'' of the title. Three further mechanisms shape the texture of the result:
neighbour support, which gives domains surface tension; a rare-type advantage within a
species, which keeps its sub-species comparable in size and so makes fragmentation a real
division; and a colour thermostat, which makes hue dynamical and heritable. Five of the model's rules are not fixed but switchable, and each
alternative is set out here beside the default.
\end{abstract}

\section{Overview}

The lattice starts with a single occupied cell of a single species at its centre.
Occupied cells spread into empty ground and contest each other's territory. Two
mechanisms generate novelty. \emph{Mutation} converts one cell into a new sub-species of
its own species, adding to that species' internal diversity. \emph{Fragmentation}
converts an entire species' worth of accumulated diversity into new species at a
stroke, when the species runs out of energy. Between these two mechanisms the single
initial species can, in principle, give rise to indefinitely many; in practice it gives
rise to tens of thousands in a few hundred sweeps.

Three further mechanisms shape what happens between those events. \emph{Neighbour
support} rewards a cell for being surrounded by its own kind, which gives domains
surface tension and keeps them compact. An \emph{occupancy bias} within a species
favours whichever sub-species is locally rarer, which keeps internal diversity alive
long enough to matter. And a \emph{colour thermostat} governs how the hues of new
species separate after a burst, so that the picture on screen reports the genealogy
rather than merely labelling it.

\begin{figure}[t]
\centering
\begin{minipage}[t]{0.235\linewidth}\centering
  \lattice{seq-250.png}\panelcap{$t = 250$}
\end{minipage}\hfill
\begin{minipage}[t]{0.235\linewidth}\centering
  \lattice{seq-600.png}\panelcap{$t = 600$}
\end{minipage}\hfill
\begin{minipage}[t]{0.235\linewidth}\centering
  \lattice{seq-1100.png}\panelcap{$t = 1100$}
\end{minipage}\hfill
\begin{minipage}[t]{0.235\linewidth}\centering
  \lattice{seq-3000.png}\panelcap{$t = 3000$}
\end{minipage}
\caption{A run at the default settings on a $256^2$ lattice. At $t = 250$ a single species
holds six per cent of the ground and already carries $443$ sub-species, visible as mottling
within one narrow band of hue. By $t = 600$ it has burst six times and the fragments are
still spreading into open ground --- at this colonisation rate a species reaches the
diversity that destroys it long before it reaches the walls. The lattice closes at about
$t = 1100$; by $t = 3000$ the run holds some fifty living species out of $125{,}000$ ever
created.}
\label{fig:sequence}
\end{figure}

\section{State}

\subsection{Lattice}

A square lattice of dimension $N \times N$, comprising $N^2$ cells, with $N$ a power of
two drawn from $\{64, 128, 256, 512\}$; the default is $N = 256$. The neighbourhood is
\textbf{Moore} --- the eight cells that touch, edges and corners alike --- and the
boundary is \textbf{periodic}: the lattice is a torus, with no edge and no corner. Each
cell is either \textbf{empty} or \textbf{occupied}.

\subsection{Taxonomy}

Every occupied cell carries two labels: a \textbf{species} and, within that species, a
\textbf{sub-species}. The two levels are nested --- a mutation changes a cell's
sub-species label while leaving its species label intact --- and they differ only in how
they interact (Section~\ref{sec:competition}). Sub-species spread through the domain of
their parent species by much the same mechanism by which the species spreads into empty
ground.

In the lattice itself only the sub-species label is stored; the species is recovered
from it, since each sub-species belongs to exactly one species at a time. When a species
fragments, its sub-species are re-pointed at their new parents, and every cell follows
without being touched.

\subsection{Species-level bookkeeping}

For each extant species $i$ the model tracks

\begin{center}
\begin{tabular}{ll}
\toprule
$\alpha_i$ & \emph{energy} of species $i$: a real variable, evolving in time \\
$N_i$      & number of lattice cells occupied by species $i$ \\
$X_i$      & number of \emph{extant} sub-species of species $i$ \\
$\varphi_i$ & the species' hue, in degrees on the colour circle \\
\bottomrule
\end{tabular}
\end{center}

\noindent
and, for display only, a serial \textbf{tag} that is never reused, and a count of the
sub-species it has ever carried. Each sub-species in turn knows its own hue and the
number of cells it occupies.

Note that $X_i \le N_i$, with equality when every cell of the species is its own
sub-species. The variables $N_i$ and $X_i$ are read off the lattice; $\alpha_i$ obeys its
own equation of motion (Section~\ref{sec:energy}).

\medskip
\noindent
\textbf{$X_i$ counts what is alive.} There are two things ``number of sub-species'' could
mean: those currently extant, or every one the species has ever produced. $X_i$ is the
extant count, and the difference is not cosmetic. A sub-species driven to extinction by
internal competition stops costing its parent anything, so the coarsening of
Section~\ref{sec:typeII} is a genuine regulator of dissipation, and a species can hold its
diversity --- and hence its energy --- in balance. Were the cumulative count used instead,
the cost could only ever rise, and fragmentation would become certain from the moment a
species began to diversify. The cumulative count is tracked, but only to be displayed.

\subsection{Initial condition}

One occupied cell at the centre of the grid, belonging to a single species with a single
sub-species; all other cells empty. The founder's hue is drawn uniformly from the circle,
and its energy is set to $\alpha_0$.

\section{Time}
\label{sec:time}

The time increment associated with one elementary competition event is
\begin{equation}
  \dt = \frac{1}{N^2},
  \label{eq:dt}
\end{equation}
so that one unit of model time --- one \textbf{sweep} --- corresponds to $N^2$ elementary
events, i.e. on average one attempt per lattice site.

Competition is a discrete sequence of events; mutation, energy and colour are continuous
processes running alongside it. Rather than interleave them event by event, the model
advances in \textbf{chunks}. A chunk resolves $K$ elementary contests, then advances the
clock by $K\dt$ and applies the mutation, energy and colour updates once each over that
whole interval. The chunk size is
\begin{equation}
  K = \min\bigl(\max(256,\; 4 S),\; 4096\bigr),
\end{equation}
where $S$ is the current number of living species: the passes that cost $O(S)$ are run
less often when there are more species to pay for. The splitting is the only
approximation in the model's time-stepping, and its error is bounded by the amount
$\alpha$ can move within one chunk --- at the default settings, a few parts in a
thousand.

\section{Parameters}

\begin{center}
\begin{tabular}{llll}
\toprule
Symbol & Meaning & Default & Range \\
\midrule
$N$          & linear grid dimension                     & $256$      & $64$--$512$ \\
$d$          & resistance of empty ground; $P_{\text{expand}} = 1-d$ & $0.90$ & $0$--$1$ \\
$E$          & energy supplied per occupied cell per sweep & $0.0316$ & $10^{-3}$--$10^{-0.5}$ \\
$\beta$      & cost per sub-species per sweep            & $0.288$    & $10^{-1}$--$10^{1.5}$ \\
$m$          & mutation-rate constant                    & $8.9\times10^{-5}$ & $10^{-5}$--$10^{-1}$ \\
$\delta$     & fragmentation-rate constant               & $7.1\times10^{-3}$ & $10^{-3}$--$10^{1}$ \\
$\alpha_0$   & seed and reset energy                     & $0.01$     & $10^{-2}$--$10^{2}$ \\
$\lambda$    & rare-sub-species advantage                & $1.14$     & $0$--$2$ \\
$\theta$     & neighbour-support exponent                & $3$        & $0$--$3$ \\
\midrule
\multicolumn{4}{l}{\emph{Colour, thermostat rule}} \\
$D$          & burst hue dispersal, degrees              & $90$       & $0$--$360$ \\
$T$          & colour settle time, sweeps                & $2$        & $0.32$--$200$ \\
$a$          & settle pull \ (labelled $\alpha$ in the panel) & $2$        & $0$--$12$ \\
$k$          & number of settled hues                    & $12$       & $2$--$36$ \\
\midrule
\multicolumn{4}{l}{\emph{Colour, other rules}} \\
$g$          & burst hue spread                          & $2$        & $0$--$5$ \\
$\sigma$     & burst hue jitter, degrees                 & $47$       & $0$--$180$ \\
\bottomrule
\end{tabular}
\end{center}

\noindent
All rates are per unit of model time, i.e. per sweep. The sliders for $E$, $\beta$, $m$,
$\delta$, $\alpha_0$ and $T$ travel on a logarithmic scale, since each matters across
decades rather than across an interval.

\section{Competition: the elementary update}
\label{sec:competition}

At each iteration a cell is chosen uniformly at random from all $N^2$, and one of its
eight neighbours is chosen uniformly at random. If the two carry the same label, or both
are empty, nothing happens. Otherwise the two cells \textbf{compete}, and the label of
the winner is copied into the cell of the loser. Which of the two wins is probabilistic.
There are three cases.

\paragraph{Type I: empty--occupied.} The occupied cell takes the empty one with a fixed
probability, independent of everything else. Empty ground behaves as a rival of fixed
strength $d$ --- its \emph{resistance} --- so that
\begin{equation}
  P_{\text{expand}} = 1 - d,
  \qquad
  P_{\text{no expand}} = d.
  \label{eq:typeI}
\end{equation}
At the default $d = 0.90$ a contested empty cell is taken one time in ten, which makes
colonisation slow: the lattice takes some eleven hundred sweeps to fill.
Colonisation carries no dependence on energy. This is a deliberate restriction: $\alpha$
governs conflict between species and nothing else, so growth into open ground proceeds at a
uniform rate and a species cannot widen its frontier by being rich. It can only take ground
that is already held.

The asymmetry of this case must be stated explicitly: if the occupied cell loses, the
outcome is simply that \emph{no expansion occurs}. The occupied cell is not replaced by an
empty one; empty ground never invades. Note also that either of the two cells may be the
focal one --- an empty cell chosen first is colonised by an occupied neighbour just as
readily --- so the rule is symmetric in the pair even though its outcome is not.

\paragraph{Type II: two sub-species of the same species.}
\label{sec:typeII}
Two sub-species of a common species share their parent's energy, so $\alpha$ cannot
separate them. Stripped of everything else the contest is therefore a fair coin,
\begin{equation}
  P_{u \text{ wins}} = P_{v \text{ wins}} = \tfrac12,
  \label{eq:typeIItie}
\end{equation}
under which competition within a species is unbiased and sub-species boundaries perform a
neutral, voter-model coarsening inside the species' domain.

Two factors may be layered on top of that coin, each with its own control: how well each
cell is supported by its immediate neighbours, and how rare each sub-species is within the
parent. With both in force --- the default --- and writing $u$ and $v$ for the two
sub-species, $n_u$ and $n_v$ for the number of Moore neighbours matching each at its own
cell, and $|u|$, $|v|$ for the number of cells each occupies,
\begin{equation}
  P_{u \text{ wins}}
  = \frac{(n_u+1)^{\theta}\, e^{-\lambda |u| / \kappa_i}}
         {(n_u+1)^{\theta}\, e^{-\lambda |u| / \kappa_i}
        + (n_v+1)^{\theta}\, e^{-\lambda |v| / \kappa_i}},
  \qquad
  \kappa_i = \frac{N_i}{X_i}.
  \label{eq:typeII}
\end{equation}
The parent's energy has cancelled from numerator and denominator, as it must. The two
factors are independent and are switched separately: the \emph{neighbour support} control
removes the first, which is equivalent to setting $\theta = 0$, and the
\emph{within-species} control's \emph{tie} setting removes the second, which is equivalent
to $\lambda = 0$. Only when both are removed is \eqref{eq:typeIItie} recovered exactly ---
support alone still favours the better-embedded cell, even under the tie setting.

The exponential factor is a \textbf{rare-type advantage}. Because $|u|$ appears with a
negative sign, the \emph{larger} sub-species is the one penalised: at the default
$\lambda = 1.14$, a sub-species that exceeds its rival by one mean sub-species area wins
only about a quarter of their contests. The normalisation $\kappa_i$ is the mean sub-species area
within the parent, so the bias reads differences in size relative to the species' own
internal scale rather than in absolute cells.

What the advantage changes is not how many sub-species a species carries but how large any
one of them is allowed to grow. Under the tie setting, coarsening within a species does what
a voter model does: one sub-species engulfs the domain, and most of the rest are recent
mutations that have not yet been swallowed. At $t = 2000$ and the default parameters a single
species has then taken the whole lattice, $88$ per cent of it is one sub-species, and $96$
per cent of its sub-species are lone cells. Under the advantage the largest species is
instead a mosaic of comparable patches --- median area ten cells, largest twenty-nine ---
because nothing can grow large without being penalised for it. Section~\ref{sec:consequences} draws out why this matters:
a burst divides a species along its sub-species, so it is their size distribution, and not
their number, that decides how much of the parent survives the catastrophe.

Whichever settings are in force, this is the sole channel by which a sub-species can be
driven extinct without the whole species being involved.

\paragraph{Type III: two distinct species.}
For species $A$ and $B$,
\begin{equation}
  P_{A \text{ wins}} = \frac{\alpha_A (n_A+1)^{\theta}}
                            {\alpha_A (n_A+1)^{\theta} + \alpha_B (n_B+1)^{\theta}},
  \label{eq:typeIII}
\end{equation}
where $n_A$ and $n_B$ now count Moore neighbours of the \emph{same species} at each of
the two cells. A species with larger $\alpha$ is more ``powerful'': it conquers its
neighbours with greater likelihood. This is the sense in which $\alpha$ behaves like an
energy. Should both strengths vanish, the contest is decided by a fair coin.

The occupancy bias does \emph{not} reach this case. Rarity is an advantage inside a
species, where it preserves the diversity that drives the model; between species it would
merely blunt the effect of $\alpha$, which is the quantity the model is about.

\subsection{Neighbour support}

Both Types II and III carry the factor $(n+1)^{\theta}$, where $n \in \{0,\dots,8\}$
counts the neighbours of the contesting cell that match it at the level of that contest
--- sub-species for Type II, species for Type III. Empty neighbours are not counted, so a
cell on the frontier is weaker than one in the interior. The support factor therefore
ranges over $[1, 9^{\theta}]$, and at the default $\theta = 3$ it spans nearly three orders of
magnitude: a cell with all eight neighbours of its own kind beats an isolated intruder
in better than 998 contests out of 1000.

The effect is surface tension. Without it, boundaries wander freely and domains dissolve
into filigree; with it, a domain pays a price along its perimeter and none in its
interior, so it rounds up, holds its shape, and can only be eaten from the edge. Setting
$\theta = 0$ switches it off, and Types II and III revert to their unadorned forms.

\begin{figure}[t]
\centering
\begin{minipage}[t]{0.47\linewidth}\centering
  \lattice{support-on.png}\panelcap{support on, $\theta = 3$}
\end{minipage}\hfill
\begin{minipage}[t]{0.47\linewidth}\centering
  \lattice{support-off.png}\panelcap{support off, $\theta = 0$}
\end{minipage}
\caption{The same seed and the same time ($t = 2000$), with neighbour support on and off.
With it, domains are compact and their boundaries short: seven per cent of occupied cells
lie on a species boundary. Without it, boundaries wander, domains interpenetrate, and the
boundary fraction rises to thirty-six per cent. The ragged state also carries far more
species at once --- $1335$ against $26$ --- because nothing gathers a domain together and
small holdings are never squeezed out.}
\label{fig:support}
\end{figure}

\subsection{A unified reading of the three rules}

Types II and III are one rule. Give every occupied cell $c$, carrying sub-species $u$ of
species $i$, a strength
\begin{equation}
  w(c) = \underbrace{\alpha_i}_{\text{species energy}}
       \times \underbrace{(n_c+1)^{\theta}}_{\text{local support}}
       \times \underbrace{e^{-\lambda |u| / \kappa_i}}_{\text{rare-type discount}},
\end{equation}
and let the contest go to $c$ with probability $w(c) / (w(c) + w(c'))$. Then:

\begin{center}
\begin{tabular}{lll}
\toprule
Opponent & What survives in $w$ & Recovers \\
\midrule
a sub-species of the same species & support and rarity; $\alpha_i$ cancels & \eqref{eq:typeII} \\
a distinct species $B$            & energy and support; rarity omitted     & \eqref{eq:typeIII} \\
\bottomrule
\end{tabular}
\end{center}

\noindent
Type I stands outside this scheme. It is not a contest between strengths but a fixed
coin of bias $d$, and the loser of it forfeits nothing, because empty ground has nothing
to lose.

\section{Mutation: the origin of sub-species}
\label{sec:mutation}

A sub-species arises when a single cell mutates within the occupied domain of a species.
The mutation changes that cell to a new sub-species type while it retains its primary
species classification, and $X_i \mapsto X_i + 1$. The new sub-species is given a hue one
small step from its parent's --- uniformly within $\pm 7^{\circ}$ --- and thereafter
spreads by Type II competition.

The rate may be set in either of two ways, chosen by the \emph{mutation rate} control:
\begin{enumerate}[label=(\alph*),itemsep=2pt]
  \item a \textbf{constant rate per cell}, giving a species-wide rate $m N_i$;
  \item a \textbf{rate per cell proportional to energy}, $m\,\alpha_i$ per cell, giving a
        species-wide rate
        \begin{equation}
          R^{\text{mut}}_i = m\,\alpha_i N_i .
        \end{equation}
\end{enumerate}
Option (b) is the default. It couples novelty to prosperity, so that rich species
diversify faster and thereby hasten their own dissipation, which is the feedback the
model exists to display.

\paragraph{Implementation.} Rather than tracking a mutation clock for every cell,
mutation events are drawn against the population-wide rate and then assigned. An
exponential clock accumulates: each chunk subtracts $R\,\dt$ from a running debt, and
whenever the debt falls to zero an event fires and a fresh $\mathrm{Exp}(1)$ draw is added
back. This is exact for a rate held constant across the chunk.

When an event fires, a cell is chosen uniformly at random from all occupied cells. Under
option (a) this alone gives each species its share, since a species is selected in
proportion to its area. Under option (b) the global rate is run at the envelope
$m\,\alpha_{\max} N_{\text{occ}}$, where $\alpha_{\max}$ is the largest energy currently
held by any species, and the drawn event is accepted with probability
$\alpha_i / \alpha_{\max}$. This \emph{thinning} recovers the per-cell rate $m\alpha_i$
exactly, at the cost of some rejected draws.

One consequence deserves note: if the mutating cell was the last of its sub-species, that
sub-species dies as the new one is born, and $X_i$ is unchanged. Mutation adds diversity
only where there is something to mutate away from.

\section{Energy}
\label{sec:energy}

Each occupied cell supplies energy at rate $E$, so a species of area $N_i$ has a total
energy capacity per unit time of $N_i E$. Each extant sub-species costs $\beta$ per unit
time --- the \emph{cost of multiplicity}. Hence
\begin{equation}
  \frac{d\alpha_i}{dt} = N_i E - \beta X_i .
  \label{eq:ode}
\end{equation}
Discretising \eqref{eq:ode} over a chunk of duration $\dt_K = K/N^2$,
\begin{equation}
  \alpha_i(t + \dt_K) \;=\; \max\Bigl(0,\;
    \alpha_i(t) + \dt_K \bigl( N_i E - \beta X_i(t) \bigr)\Bigr).
  \label{eq:euler}
\end{equation}

\paragraph{Remarks on the energy equation.}
\begin{enumerate}[label=(\roman*),itemsep=3pt]
  \item $\alpha_i$ does not appear on the right-hand side of \eqref{eq:ode}. It is a pure
        accumulator --- a running store of energy --- not a self-regulating quantity, and
        it has no fixed point of its own. It grows while $X_i < N_i E / \beta$ and falls
        once diversity passes that threshold.
  \item Since $X_i \le N_i$, the most negative attainable drift is
        \begin{equation}
          \left.\frac{d\alpha_i}{dt}\right|_{\text{extreme}} = N_i\,(E - \beta),
          \qquad \text{attained at } X_i = N_i .
        \end{equation}
        It follows that a species can lose energy only if $\beta > E$. This is the
        condition under which the catastrophe of Section~\ref{sec:fragment} is reachable
        at all, and the interface warns when it is violated.
  \item A species carrying a single sub-species needs at least $\beta/E$ cells to break
        even. Below that area its energy drains away whatever it does. At the default
        settings the figure is about nine cells, and it governs the fate of the smallest
        fragments of a burst (Section~\ref{sec:consequences}).
  \item Energy is floored at zero rather than allowed to go negative. A species at
        $\alpha_i = 0$ loses every Type III contest it enters against any species with
        energy to spend, and under the $\alpha$-coupled mutation rule it stops mutating
        altogether. It is not dead --- it still takes empty ground, and if its area
        exceeds $\beta X_i / E$ its energy will recover --- but it is defenceless until
        it does.
  \item Area helps and diversity hurts, and both linearly. A species survives by being
        large and internally uniform.
\end{enumerate}

\section{Fragmentation}
\label{sec:fragment}

When a species runs out of energy --- because it has come to carry too many sub-species ---
it \textbf{fragments} in a catastrophe event. Every sub-species immediately breaks away and
becomes a new, independent species. The domains do not move: what changes is the law by
which they interact, which switches from Type II to Type III. Neighbouring patches that
a moment earlier were contesting a shared inheritance now contest each other in
proportion to their separate energies.


\begin{figure}[t]
\centering
\begin{minipage}[t]{0.47\linewidth}\centering
  \lattice{burst-before.png}\panelcap{before}
\end{minipage}\hfill
\begin{minipage}[t]{0.47\linewidth}\centering
  \lattice{burst-after.png}\panelcap{after}
\end{minipage}
\caption{A burst, across the single animation frame in which it happens, at
$t \approx 1239$. The species held $12{,}826$ cells --- a fifth of the lattice --- and still
had $\alpha = 755$ in hand when the hazard \eqref{eq:fragrate} fired; all $1444$ of its
sub-species became species at once, and it was the only species to fragment in that frame.
Nothing has moved: the outline of the parent is exactly preserved, and its neighbours are
untouched. What has changed is the law inside that outline, which a moment earlier was
Type II and is now Type III. This panel is drawn under the redraw colour rule rather than the
default, so that each fragment takes a fresh hue and the division is unmistakable; under the
default the fragments would scatter about their parent's own hue instead.}
\label{fig:burst}
\end{figure}

\subsection{The triggers}

There are two triggers. They are not alternatives: both are live at once, and they are
tested in the order given.

\begin{enumerate}[label=(\alph*),itemsep=3pt]
  \item \textbf{Deterministic.} Fragment as soon as the update \eqref{eq:euler} would
        take the energy to zero or below:
        \begin{equation}
          \alpha_i(t) + \dt_K \bigl( N_i E - \beta X_i \bigr) \le 0
          \;\Longrightarrow\; \text{burst.}
        \end{equation}
  \item \textbf{Stochastic.} Otherwise fragment with probability
        \begin{equation}
          1 - e^{-R^{\text{frag}}_i \dt_K},
          \qquad
          R^{\text{frag}}_i \;=\; \delta\,\frac{X_i}{\alpha_i(t+\dt_K)},
          \label{eq:fragrate}
        \end{equation}
        a rate increasing in diversity and decreasing in energy, evaluated at the newly
        updated energy.
\end{enumerate}

\noindent
Rule (b) diverges as $\alpha_i \to 0^{+}$, so it anticipates rule (a) rather than
competing with it, while also permitting early bursts in a species that is merely
strained. At the default settings rule (b) does essentially all the work: a species in
the steady regime of Section~\ref{sec:consequences} holds an energy of several hundred
and never approaches zero, so it bursts because \eqref{eq:fragrate} eventually fires, not
because it is bankrupt. Rule (a) is the floor beneath that, and matters chiefly for
species too small to pay for the diversity they carry.

\paragraph{A species with one sub-species cannot fragment.} There must be at least two
sub-species for a burst to mean anything, and the trigger is suppressed below that. A
monomorphic species that exhausts its energy simply sits at $\alpha_i = 0$ until its area
grows enough to recover, or until it is overrun.

\subsection{What the fragments inherit}

Each new species takes exactly the cells its sub-species held, is given a fresh serial
tag, and begins life with $X = 1$: it is its own only sub-species. Its hue is set by the
colour rule of Section~\ref{sec:colour}. Its energy is the one thing the burst does not
settle by itself, since the parent's store is being wound up rather than divided by any
natural rule; three readings are offered, chosen by the \emph{energy of the fragments}
control:

\begin{center}
\begin{tabular}{lll}
\toprule
Rule & Each fragment receives & \\
\midrule
reset      & $\alpha_0$                                     & (default) \\
equal split& $\alpha_{\text{parent}} / X_{\text{parent}}$    & \\
by area    & $\alpha_{\text{parent}} \cdot |u| / N_{\text{parent}}$ & \\
\bottomrule
\end{tabular}
\end{center}

\noindent
The two splitting rules interact with the choice of trigger in a way worth noticing.
A deterministic burst happens precisely when the parent's energy has reached zero, so
there is nothing left to divide and both splitting rules hand every fragment an energy of
zero. Only a stochastic burst, which fires while the parent is still solvent, leaves an
inheritance. The reset rule is therefore the only one that behaves identically under both
triggers, which is why it is the default.

Once every sub-species has been promoted, the parent is emptied of cells and sub-species
and removed from the population. The time of the burst is logged, so that it can be
drawn on the graph.

\section{Colour}
\label{sec:colour}

Cells are coloured in HSB with saturation fixed at $0.72$ and brightness at $0.95$, and
hue free, so that the entire taxonomy is carried on the hue circle: red at $0^\circ$,
through orange, yellow, green, blue and indigo, to violet and back to red at $360^\circ$.
Empty ground is drawn near-black.

Within a species the arrangement is fixed. The colour of a new
sub-species is a small step in hue from the antecedent in which the mutation occurred ---
uniformly within $\pm 7^{\circ}$ --- so hue measures genealogical distance from the
founder and the species reads on screen as a band of closely related shades.

What happens \emph{at} a burst is a matter of choice, and four rules are offered, selected
by the \emph{colour after a burst} control. Under each of them, the sub-species that is
about to become a species holds the hue it has until the moment it is promoted.

\begin{description}[itemsep=4pt,leftmargin=1.4em]
\item[Redraw.] Each fragment draws a fresh hue uniformly from the circle and keeps it.
  The taxonomy is legible at a glance --- no two fragments will be confused --- but the
  genealogy is discarded at every burst.

\item[Jitter.] Each fragment steps off its own hue by a normal draw of standard deviation
  $\sigma$. The family stays recognisable; the siblings separate. This is the default, and
  the cost of it is that neighbouring species are often close in hue --- which is what the
  perimeter of Section~\ref{sec:interface} is for.

\item[Spread.] Each fragment is pushed $g$ times its own offset from the family's mean
  hue, in the direction it already lay:
  \begin{equation}
    \varphi \;\mapsto\; \bar\varphi + g\,\bigl(\varphi - \bar\varphi\bigr),
  \end{equation}
  with the difference taken as the signed separation in $(-180^{\circ}, 180^{\circ}]$ and
  $\bar\varphi$ the \emph{circular} mean --- the direction of the resultant of the unit
  vectors, since $350^{\circ}$ and $10^{\circ}$ average to $0^{\circ}$ and not to
  $180^{\circ}$. The siblings fan apart along the lines their mutations had already
  started, so the branching structure survives the catastrophe that produced it.

\item[Thermostat.] The fragment keeps its hue and inherits, instead, a quantity of
  \emph{heat} that its parent's multiplicity paid for. It then walks away from its
  siblings as that heat decays. Alone among the four it makes hue a dynamical variable in
  its own right, obeying a stochastic differential equation on the colour circle; it is
  described below.
\end{description}


\begin{figure}[t]
\centering
\begin{minipage}[t]{0.235\linewidth}\centering
  \lattice{hue-redraw.png}\panelcap{redraw}
\end{minipage}\hfill
\begin{minipage}[t]{0.235\linewidth}\centering
  \lattice{hue-thermo.png}\panelcap{thermostat}
\end{minipage}\hfill
\begin{minipage}[t]{0.235\linewidth}\centering
  \lattice{hue-spread.png}\panelcap{spread}
\end{minipage}\hfill
\begin{minipage}[t]{0.235\linewidth}\centering
  \lattice{hue-jitter.png}\panelcap{jitter}
\end{minipage}
\caption{The four colour rules, same seed, same time ($t = 2000$). The patterns are not
merely similar but identical --- the lattices agree cell for cell, as the two-stream
construction of Section~\ref{sec:numerics} requires --- so the panels differ in hue alone.
Redraw scatters fragments across the whole circle. The thermostat holds each species on one
of $k = 12$ settled hues, giving a quantised palette in which relatives share a colour.
Spread fans siblings apart from their family mean, and jitter --- the default --- steps each
off its own hue, so both keep some memory of descent.}
\label{fig:colour}
\end{figure}

\subsection{The thermostat}

Under this rule each species' hue $\varphi$ obeys, in degrees and with time in sweeps,
\begin{equation}
  d\varphi \;=\; -\,a \sin\!\left(k \varphi \frac{\pi}{180}\right) dt
              \;+\; \sqrt{r(t)}\;dW .
  \label{eq:huesde}
\end{equation}
The two terms do different jobs.

\paragraph{The noise, and the heat that drives it.} A fragment of a parent that carried
$X$ sub-species is given a total variance budget
\begin{equation}
  V(X) = D^2 \frac{X}{8},
\end{equation}
so that its total r.m.s. displacement in hue is $D\sqrt{X/8}$ degrees. The reference
multiplicity of $8$ makes $D$ the plain design dial: a species that bursts into eight
fragments disperses them by $D$ degrees, one that bursts into thirty-two disperses them
by twice that. Multiplicity therefore buys colour separation in exactly the way it buys
catastrophe.

That budget is spent over a settling window of $T$ sweeps, at a variance rate falling
linearly to nothing:
\begin{equation}
  r(t) = \frac{2V}{T}\left(1 - \frac{t}{T}\right), \qquad 0 \le t \le T,
\end{equation}
whose integral over the window is exactly $V$. When the window closes the motion has
eased to a stop rather than being cut off, and the heat is set to exactly zero: a settled
species holds its colour precisely until it fragments in turn.

\paragraph{The drift, and the wells.} Left to itself the noise would stop wherever it
happened to be. The drift term gathers it instead. The function $-a\sin(k\varphi)$ has
stable zeros at the $k$ hues evenly spaced around the circle, at multiples of
$360^{\circ}/k$, and pulls the walk towards whichever is nearest with a strength set by
$a$. A fragment does not merely stop diffusing; it is drawn into a definite colour and
held there. When the settling window closes the hue is placed exactly on its well, so the
palette stays crisp rather than merely approximate. Setting $a = 0$ removes the drift and
leaves the walk to stop where it lies.

\paragraph{The band moves as a body.} The walk is integrated once per species, and the
resulting increment is added to the species' hue \emph{and} to each of its sub-species'
hues alike. The $\pm 7^{\circ}$ fine structure that mutation builds inside a species
therefore survives the walk intact: the band drifts and settles as a unit, and its
internal genealogy is not smeared away.

\medskip
\noindent
\textbf{A road not taken.} One could go further and dispense with discrete labels
altogether, defining the model \emph{by intervals of hue}: species identity would then be a
matter of proximity on the colour circle, and speciation a matter of a lineage's hue
drifting far enough from its relatives to fall outside the interval. The model does not do
this, and the distinction is worth stating plainly. Even under the thermostat, hue is a
consequence of the taxonomy and never a cause: it is carried along by species that the
competition rules have already made, and no cell ever changes hands because of a
colour.

\section{The algorithm}
\label{sec:algorithm}

The run is a sequence of chunks. For each chunk, with $K$ as in Section~\ref{sec:time}:

\begin{enumerate}[itemsep=4pt]
  \item \textbf{Competition.} Repeat $K$ times: choose a cell uniformly at random, and
        one of its eight neighbours uniformly at random. If the labels agree, or both are
        empty, do nothing. Otherwise classify the pair as Type I, II or III, resolve the
        contest with the appropriate probability, and copy the winner's labels into the
        loser's cell. Update the affected $N_i$, and $X_i$ if a sub-species has just been
        driven extinct; retire any species that has lost its last cell.
  \item \textbf{Advance the clock} by $\dt_K = K/N^2$.
  \item \textbf{Mutation.} Draw events against the population-wide rate over $\dt_K$. For
        each, pick a uniformly random occupied cell, apply the thinning test if the rate
        is $\alpha$-coupled, and on acceptance give the cell a fresh sub-species label and
        a hue one small step from its parent's.
  \item \textbf{Energy.} Update every $\alpha_i$ by \eqref{eq:euler}, and record the
        largest for the next chunk's thinning envelope. Test each species for
        fragmentation by the two triggers of Section~\ref{sec:fragment}, collecting those
        that fire; then fragment them, promoting each sub-species to an independent
        species.
  \item \textbf{Colour.} Integrate \eqref{eq:huesde} over $\dt_K$ for every species that
        still has heat left, subdividing the step if the drift demands it.
  \item \textbf{Sample.} Whenever the clock passes a multiple of two sweeps, record a
        snapshot for the graph.
\end{enumerate}

\noindent
The visible image is repainted cell by cell as labels change, and in full once per frame
whenever a burst or a hue movement has altered colours that were already on screen.

\section{Structural consequences}
\label{sec:consequences}

Four consequences follow from the rules as stated, and are worth keeping in view when
reading simulation output.

\paragraph{The occupied region never retreats.} No rule ever vacates a cell. Type I
cannot, by construction; Types II and III only relabel. The occupied set is therefore
monotonically non-decreasing, and the lattice fills --- at the default settings, by about
$t = 1100$. Once it is full, Type I events cease and the dynamics reduce to relabelling
under Types II and III --- a biased voter model on a full lattice, punctuated by mutation
and fragmentation.

\paragraph{The characteristic cycle.} A species grows, its area drives $\alpha$ up,
mutation accumulates sub-species, the cost $\beta X_i$ overtakes the supply $N_i E$, and
the species bursts into as many new species as it had sub-species. Those fragments, now
mutually hostile and individually small, contest the territory their parent held. The
model's macroscopic behaviour is the statistics of this cycle repeated across a lattice.
A single burst is a large event: the species of Figure~\ref{fig:burst} holds a fifth of the
lattice and promotes $1444$ sub-species at one stroke. At the default colonisation rate the
first turn of the cycle is also complete long before the lattice is full, so growth and
fragmentation overlap: the founder of Figure~\ref{fig:cycle} bursts holding a seventh of the
ground, with open country still ahead of it.


\begin{figure}[t]
\centering
\begin{tikzpicture}
\begin{groupplot}[
    group style={group size=1 by 2, vertical sep=7mm, x descriptions at=edge bottom},
    width=0.95\linewidth, height=48mm,
    xmin=0, xmax=360,
    tick align=outside, tick pos=left,
    every axis plot/.append style={line width=0.9pt},
    label style={font=\small}, tick label style={font=\footnotesize},
    axis line style={black!55},
]
\nextgroupplot[ylabel={energy $\alpha_i$}, ymin=0, ymax=850,
               ytick={0,200,400,600,800}]
  \addplot[color=blue!55!black] table[x=t, y=alpha] {figures/founder.dat};
\nextgroupplot[xlabel={time $t$ \ (sweeps)}, ylabel={sub-species}, ymin=0, ymax=1150,
               ytick={0,250,500,750,1000},
               legend style={font=\footnotesize, draw=none, fill=none,
                             at={(0.03,0.95)}, anchor=north west, legend columns=1}]
  \addplot[color=black!45, dashed] table[x=t, y=balance] {figures/founder.dat};
  \addlegendentry{$N_i E / \beta$}
  \addplot[color=red!65!black] table[x=t, y=X] {figures/founder.dat};
  \addlegendentry{$X_i$}
\end{groupplot}
\end{tikzpicture}
\caption{The founder species, followed from the first cell to its burst at $t = 356$. Its
energy climbs steeply while diversity is still cheap, peaking near $754$ at $t \approx 290$;
thereafter mutation has brought $X_i$ up to the balance point $N_i E/\beta$, the two curves
in the lower panel lock together, and $\alpha_i$ goes flat. The species is not bankrupt when
it bursts --- it still holds $\alpha = 717$ --- but by then it carries $1034$ sub-species,
and the hazard \eqref{eq:fragrate} has become large enough to fire. It is also still
growing, on only $8763$ cells --- a seventh of the lattice. The lower panel is the whole
argument of this section in one picture: the dashed line is area, rescaled, and diversity
tracks it.}
\label{fig:cycle}
\end{figure}

\paragraph{The population sits at the balance point.} The cycle does not run away. Because
mutation raises $X_i$ while nothing lowers $N_i E$, a species is driven towards the
locus
\begin{equation}
  X_i \;\approx\; \frac{N_i E}{\beta},
\end{equation}
at which production and dissipation cancel and $\alpha_i$ goes flat. Every substantial
species in a mature run is found there, to within a few per cent, with its energy parked
at a few hundred and its drift near zero. The system settles into a marginal state rather
than an oscillation, and what ends a species is not bankruptcy but the accumulating
hazard of \eqref{eq:fragrate}, which grows with the diversity the balance point obliges
it to carry.

\paragraph{Fragments are filtered by area.} A burst is not a fair division. Each fragment
inherits its sub-species' cells, and those areas are broadly distributed; but by remark
(iii) of Section~\ref{sec:energy}, a fragment of fewer than $\beta/E$ cells cannot pay for
even its own single sub-species. Its energy drains to zero within a fraction of a sweep,
after which it loses every contest it enters and is overrun. Immediately after a large
burst, hundreds of species may sit at $\alpha = 0$ simultaneously. The great majority of
the species a burst creates are gone almost at once, and only the larger fragments
survive to grow, accumulate diversity, and burst in their turn --- which is why a run
that has produced tens of thousands of species may hold only a few hundred alive.

\section{Numerical and implementation notes}
\label{sec:numerics}

\paragraph{Random numbers, and two streams.} The generator is the Park--Miller minimal
standard, $x \mapsto 16807\,x \bmod (2^{31}-1)$, evaluated by Schrage's trick to stay
within 32-bit arithmetic. It is run as \emph{two} independent streams, seeded from a
common value: one drives the dynamics, the other drives hue noise alone. The separation
is deliberate. Switching the colour rule, or the thermostat's parameters, cannot move a
single cell: for a given seed the spatial trajectory is identical under all four colour
rules. To keep that guarantee, the uniform hue draw used by the redraw rule is taken from
the dynamics stream at every burst whether or not the rule is in force, so the stream
stays aligned.

\paragraph{Exact integrals where they are available.} Three quantities that could have
been linearised are not. The mutation clock uses exponential inter-arrival times rather
than a per-step Bernoulli trial. The fragmentation test uses $1 - e^{-R\dt}$ rather than
$R\dt$, which matters when $R\dt$ is not small. And the thermostat's variance, whose rate
is linear in time, is integrated in closed form over each step as
$r_0 \tau - \tfrac12 c\,\tau^2$, clipped at the moment the settling window closes. The
total dispersal is therefore exactly $D\sqrt{X/8}$ whatever the chunk size happens to be.

\paragraph{Sub-stepping the colour drift.} The drift in \eqref{eq:huesde} stiffens as
$ak$ grows, and Euler--Maruyama is only conditionally stable. Each chunk is therefore
subdivided into steps no longer than $0.24/(a k \pi/180 + 1)$, to a maximum of $64$
sub-steps, before the walk is integrated. The bound is inactive when $a = 0$.

\paragraph{Thinning against a lagged envelope.} The mutation pass runs before the energy
pass within a chunk, so the envelope $\alpha_{\max}$ it thins against is the value from
the previous chunk. The lag is bounded by the change in $\alpha$ across one chunk,
which at the default settings is a fraction of a per cent of $\alpha_{\max}$.

\paragraph{Recycled records and stable identities.} Species and sub-species records are
pooled: a dead record is returned to a free list and reused, so memory tracks the
high-water mark of concurrent species rather than the number ever created. Because slots
are reused, every species also carries a \textbf{tag} --- a serial number that is never
reused --- and it is the tag that labels a species on screen and keys its trace on the
graph. A graph line can therefore never jump from a dead species to the live one that
inherited its storage. For the same reason each sample carries the colours of the species it
recorded rather than a reference to them: by the time a ghost trace is drawn its owner may be
long dead and its record reissued, so the palette held in the sample is the only surviving
record of what that line looked like.

\paragraph{Fixed constants.}
\begin{center}
\begin{tabular}{lll}
\toprule
Quantity & Value & Where \\
\midrule
hue step at mutation            & $\pm 7^{\circ}$, uniform & \S\ref{sec:mutation} \\
chunk size                      & $256$ to $4096$ events   & \S\ref{sec:time} \\
reference multiplicity          & $8$ sub-species          & \S\ref{sec:colour} \\
Euler--Maruyama safety factor   & $0.24$                   & \S\ref{sec:numerics} \\
maximum sub-steps per chunk     & $64$                     & \S\ref{sec:numerics} \\
graph sampling interval         & $2$ sweeps               & \S\ref{sec:interface} \\
species retained per sample     & $128$ largest            & \S\ref{sec:interface} \\
graph history                   & $300$ samples            & \S\ref{sec:interface} \\
burst log                       & $512$ most recent        & \S\ref{sec:interface} \\
saturation, brightness          & $0.72$, $0.95$           & \S\ref{sec:colour} \\
perimeter suppressed below       & $2$ device pixels a cell & \S\ref{sec:interface} \\
\bottomrule
\end{tabular}
\end{center}

\paragraph{The figures.} Every figure in this document was produced by the page's own
simulation core, driven headless rather than screen-captured, so the runs shown are the runs
the model performs. All use seed $20260903$ on a $256^2$ lattice at the default settings,
varying only the rule under discussion, and all are stepped in whole frames of $80{,}000$
events --- the page's own default speed --- so that the caveat of
Section~\ref{sec:interface} is satisfied and the page at \code{?seed=20260903} follows the
same trajectory. The lattice images are the model's own pixel buffer, enlarged without
interpolation, and the perimeter in Figure~\ref{fig:outline} is stroked by the page's own
code rather than a reconstruction of it.

One departure from the defaults should be declared. Figures~\ref{fig:sequence}
to~\ref{fig:cycle} are coloured under the \emph{redraw} rule rather than the default jitter,
because redraw gives every species a hue of its own and so makes the mechanism under
discussion legible; jitter keeps relatives close together, which is a virtue on screen and a
nuisance in a diagram. Figure~\ref{fig:colour} shows all four rules side by side, and
Figure~\ref{fig:outline} uses the defaults untouched.

\section{The interface}
\label{sec:interface}

\paragraph{Running.} \textsc{Start} and \textsc{Reset}, or the space bar and \code{r}. A
speed slider sets the number of elementary events attempted per animation frame, from
$10^3$ to $5\times10^5$, defaulting to $80{,}000$ --- something over seventy sweeps a
second on a $256^2$ lattice.

\paragraph{Choosing rules.} Five segmented controls select the model's switchable rules,
gathered for reference in Section~\ref{sec:options}:

\begin{center}
\begin{tabular}{lll}
\toprule
Control & Options & Default \\
\midrule
Mutation rate                 & $mN$ \textbar{} $m\alpha N$ & $m\alpha N$ \\
Within-species contests       & tie $\tfrac12$ \textbar{} occupancy $\lambda$ & occupancy $\lambda$ \\
Neighbour support, Types II and III & off \textbar{} on, $\theta$ & on \\
Colour after a burst          & redraw \textbar{} thermostat \textbar{} spread \textbar{} jitter & redraw \\
Energy of the fragments       & reset to $\alpha_0$ \textbar{} equal split \textbar{} by area & reset \\
\bottomrule
\end{tabular}
\end{center}

\noindent
The parameters belonging to a rule appear directly beneath it and only while that rule is
selected, so the panel shows what is currently in force and nothing else.

\paragraph{Reading the run.} The graph plots one point every two sweeps for the eight
largest living species, coloured as they are on the lattice, with vertical marks at every
burst. Beneath those, and much fainter, run the \emph{ghosts}: every other species in the
window, ranked by the largest value it ever reached on the metric being shown and capped at
the best thirty-two. A ghost that stops in mid-graph is a species that died there --- by
fragmenting, or by being overrun --- and one that reaches the right-hand edge is alive but
too small to have earned a full trace. Without them the graph would say only that a line had
ceased, never whether its owner had burst or merely slipped out of the leading eight. Three quantities may be shown --- area, energy $\alpha$, or sub-species per
species --- cycled with \code{g}. In area mode the scale is fixed to the whole lattice and
the count of empty cells is drawn as a dashed line; in the other two the ceiling
auto-scales, easing downwards after a burst so the trace does not jump. Beneath, a table
lists the six largest species by tag, with area, energy and $X$; below that, running
totals of events, occupancy, species and sub-species alive and ever created, contests by
type, mutations, bursts, and the frame and event rates. A warning line appears if
$\beta \le E$, since no species can then ever lose energy and nothing will ever burst.

\paragraph{The species perimeter.} A separate control, on by default, draws the boundary
between species over the grid (Figure~\ref{fig:outline}). It is a display option and not one
of the rules: the layer sits on its own canvas above the lattice, reads it, and changes
nothing in it.

The boundary is drawn a side at a time. Wherever two adjacent cells differ at the
\emph{species} level, each of the two species strokes that face on its own side of it, one
device pixel inside its own territory, in its own colour darkened to two-fifths. A contact
between species therefore carries two hairlines rather than one, and each domain is outlined
in a shade of itself, so the perimeter follows the thermostat for nothing extra --- the cast
is taken from the colour the species is already painted in, not recomputed from its hue.
Empty ground has no species and strokes nothing, so a lone species in open country draws its
own outline unanswered. Faces between two sub-species of one species carry nothing at all,
which is the point of the option: the hue structure inside a domain is left to speak for
itself.

Because each side belongs to a different species, faces cannot be merged into runs, and the
strokes are collected into one path per species and stroked in that species' colour. The
layer is keyed on the event count, so it is rebuilt only on frames that actually advanced the
lattice and a paused grid costs nothing. Below two device pixels to a cell the option is
suppressed --- a line on every face would cover the grid rather than describe it.


\begin{figure}[t]
\centering
\begin{minipage}[t]{0.315\linewidth}\centering
  \lattice{outline-off.png}\panelcap{off}
\end{minipage}\hfill
\begin{minipage}[t]{0.315\linewidth}\centering
  \lattice{outline-on.png}\panelcap{on}
\end{minipage}\hfill
\begin{minipage}[t]{0.315\linewidth}\centering
  \lattice{outline-detail.png}\panelcap{on, $44$ cells across}
\end{minipage}
\caption{The species perimeter, over the same lattice at $t = 2000$, drawn at the default
settings, with a magnified window at the right. Every species draws the boundary of its own
territory, one device pixel inside its own cells, in a cast of its own colour darkened to
two-fifths. Where two species meet you therefore see two hairlines side by side, each in its
owner's hue, and against empty ground you would see only one. The left pair shows why the
option earns its place under the default colour rule: sibling species sit close on the hue
circle, and several of the apparently single domains are two or three.
Faces between two sub-species of a single species carry nothing, which is why the mottling
inside each domain is left undisturbed: the layer reads the lattice through the same species
label a Type III contest reads, and changes nothing in it.}
\label{fig:outline}
\end{figure}

\paragraph{Typing values.} Every slider's readout is editable in place. A typed value is
mapped back onto the slider's own scale --- logarithmic for the decade parameters,
$\log_2$ for the grid size, the identity otherwise --- and is then taken exactly: neither
rounded to the slider's step nor clamped to its range, so a value can be set that the handle
could not reach. A value outside the range leaves the handle resting at the nearer end of the
track, and dragging brings it back.

\paragraph{Reproducing a run.} The whole configuration can be given in the URL:
\code{seed} fixes the random seed, which otherwise comes from the clock;
\code{size} selects the lattice ($64$, $128$, $256$ or $512$);
\code{d}, \code{E}, \code{beta}, \code{m}, \code{delta}, \code{a0}, \code{lam} and
\code{theta} set the model parameters;
\code{disp}, \code{tset}, \code{pull}, \code{wells}, \code{mult} and \code{sig} set
$D$, $T$, $a$, $k$, $g$ and $\sigma$;
and \code{rules} takes a string of digits selecting, in order, the fragment-energy rule,
the mutation rate, the within-species rule, neighbour support, and the colour rule. The
defaults correspond to \code{rules=01113}.

One caveat on reproducing a run. The seed fixes the random streams, but not the order in
which they are drawn from: the fragmentation test of Section~\ref{sec:fragment} spends one
draw per living species per chunk, and the chunk boundaries depend on how many events each
frame resolves. Two runs from one seed therefore agree exactly only if the speed slider is
left where it was. Move it and the run is still the same model, statistically, but no longer
the same trajectory.

\section{Options at a glance}
\label{sec:options}

Five of the model's rules are switchable rather than fixed. Each is described where it
belongs; this section gathers them for reference. Every figure in this document uses
these defaults except where its caption says otherwise.

\begin{center}
\begin{tabular}{llll}
\toprule
Rule & Settings & Default & Section \\
\midrule
Mutation rate & $m N_i$ \textbar{} $m \alpha_i N_i$
  & $m \alpha_i N_i$ & \S\ref{sec:mutation} \\
Within-species contests & tie ($\lambda = 0$) \textbar{} rare-type advantage $\lambda$
  & advantage & \S\ref{sec:competition} \\
Neighbour support & off ($\theta = 0$) \textbar{} on, exponent $\theta$
  & on & \S\ref{sec:competition} \\
Colour after a burst & redraw \textbar{} thermostat \textbar{} spread \textbar{} jitter
  & jitter & \S\ref{sec:colour} \\
Energy of the fragments & reset to $\alpha_0$ \textbar{} equal split \textbar{} by area
  & reset & \S\ref{sec:fragment} \\
\bottomrule
\end{tabular}
\end{center}

\noindent
Two remarks on the first three. Neighbour support applies to Types II and III alike, while
the within-species control reaches Type II only; the two are independent, and a Type II
contest is a fair coin only when both are off. And the choice of fragment energy interacts
with the fragmentation triggers: a deterministic burst leaves nothing to divide, so the two
splitting rules then hand every fragment an energy of zero, and only the reset rule behaves
identically under both triggers.

\paragraph{One condition is not optional.} Since $X_i \le N_i$, the steepest decline a
species can suffer is $N_i(E - \beta)$, so energy can fall at all only if $\beta > E$. If
$\beta \le E$, production exceeds dissipation whatever the diversity: no species ever
exhausts its store, nothing fragments, and the run reduces to the lattice filling and then
relabelling indefinitely under Types II and III. The defaults satisfy $\beta > E$ by an
order of magnitude, and the interface warns when the condition is violated. Beyond it the
model is not delicate --- the balance point of Section~\ref{sec:consequences} is an
attractor, not a tuning --- but below it there is no model at all.

\paragraph{Why the less obvious mechanisms are there.} Three of the rules above are not
forced by the model's premise, and each earns its place by what goes wrong without it. Neighbour
support gives domains surface tension; without it boundaries wander freely and the lattice
dissolves into filigree (Figure~\ref{fig:support}). The rare-type advantage $\lambda$ decides
whether a burst is a catastrophe or merely a renaming: without it one sub-species engulfs
its species, so that of the $7262$ fragments the largest species would yield at $t = 2000$
exactly one is large enough to survive, while with it about half of $1534$ fragments clear
the threshold. Diversity in the bare sense survives either way --- $X_i$ climbs to the balance
point under both settings --- but only under the advantage does fragmentation actually
multiply the species. And the colour thermostat makes hue dynamical and
heritable, so that the picture reports genealogy rather than merely labelling it --- the one
mechanism of the three that changes nothing about where a cell goes, only about what it
looks like when it gets there.

\end{document}
